\documentclass[12pt]{article}
\usepackage[margin=1in]{geometry}
\usepackage{graphicx}
\usepackage{caption}
\usepackage{hyperref}
\usepackage{float}
\usepackage{parskip}
\usepackage{booktabs}
\usepackage{array}

\hypersetup{
    colorlinks=false,
    hidelinks
}

% Renders an app screenshot if present in report/figures/, otherwise a
% labelled placeholder so the document always compiles.
\newcommand{\appshot}[3]{%
  \begin{figure}[H]
    \centering
    \IfFileExists{figures/#1}{%
      \includegraphics[width=0.82\textwidth,height=0.5\textheight,keepaspectratio]{figures/#1}%
    }{%
      \fbox{\parbox[c][4.5cm][c]{0.82\textwidth}{\centering\itshape
        [\,Screenshot placeholder\,]\\[6pt]
        Add \texttt{report/figures/#1} to display this figure.}}%
    }
    \caption{#2}
    \label{#3}
  \end{figure}%
}

\title{\textbf{GreenStep Employee Sustainability Challenge}\\[0.4em]
\large Project Report}
\author{Oli Gurmessa \and Hadi Shaar \and Naisha Srivastav \and Nathanael Agbemadon}
\date{CISC480 --- Senior Capstone\\[0.3em]
Instructor: Dr.\ Jimenez \quad\textbullet\quad Spring 2026\\[0.3em]
Submitted: May 22, 2026}

\begin{document}

\maketitle
\tableofcontents
\newpage

% ─────────────────────────────────────────────────────────────────────────────
\section{Project Overview}
% ─────────────────────────────────────────────────────────────────────────────

GreenStep Employee is a gamified, team-based web application that replaces
manual, spreadsheet-based sustainability tracking with a fast digital
platform. It was built as a CISC480 Senior Capstone project for the community
partner Minnesota GreenStep, and extends prior human-centred design work
produced by a CISC489 team~\cite{cisc489}.

\subsection{Problem}

Workplace sustainability challenges (such as Earth Month) are commonly tracked
with shared spreadsheets. Employees record their actions in a grid, and an
organizer manually tallies the results at the end of the period. This approach
suffers from ``spreadsheet fatigue'': data entry is tedious, there is no
real-time feedback, engagement drops over the course of the challenge, and
administrators find it difficult to aggregate and report on participation.

\subsection{Purpose}

The purpose of GreenStep is to replace the spreadsheet with a low-friction,
gamified digital tracker. By giving employees immediate feedback---points,
badges, streaks, and live leaderboards---the platform aims to sustain
engagement throughout a challenge while making aggregation and reporting
effortless for administrators.

\subsection{Users}

The system targets two primary user groups, consistent with the personas
defined in the prior CISC489 design work~\cite{cisc489}:

\begin{itemize}
    \item \textbf{Employees} (e.g.\ the ``Maria'' and ``Sarah'' personas) ---
    everyday participants who log sustainable actions, build streaks, and
    compete on team and individual leaderboards.
    \item \textbf{Sustainability administrators} (e.g.\ the ``David'' persona)
    --- organizers who create challenges, review and approve team requests,
    and export participation data for reporting.
\end{itemize}

Division or department leaders are a secondary audience, served indirectly
through team-based leaderboards that group participants by team.

\subsection{Constraints}

The project operated under several constraints that shaped its design:

\begin{itemize}
    \item \textbf{Zero budget} --- all infrastructure runs on free service
    tiers (Vercel, Neon, Clerk).
    \item \textbf{One semester} --- a fixed, short development timeline.
    \item \textbf{Four-person team} --- limited engineering capacity.
    \item \textbf{Self-reported data} --- actions are not independently
    verified, by partner agreement.
    \item \textbf{No dedicated DevOps} --- the solution must be low-maintenance
    and require no ongoing operational staffing.
\end{itemize}

\subsection{Proposed Solution}

GreenStep is a single team-based web application in which employees log
actions to earn points, badges, and streaks; compete on team and individual
leaderboards; and participate in one or more challenges. Administrators create
challenges, approve team-creation requests, and export participation data as
CSV. The entire system is delivered as one deployable unit, requiring no
installation by end users.

\subsection{Technology}

The application is built on a modern, type-safe web stack:

\begin{itemize}
    \item \textbf{Next.js 15} (App Router) with \textbf{TypeScript} for both
    the user interface and the API routes~\cite{nextjs}.
    \item \textbf{Tailwind CSS} for styling~\cite{tailwind}.
    \item \textbf{Clerk} for authentication and session
    management~\cite{clerk}.
    \item \textbf{Prisma} as the ORM~\cite{prisma}, backed by
    \textbf{Neon} serverless PostgreSQL~\cite{neon}.
    \item \textbf{Recharts} for the performance chart~\cite{recharts}.
    \item \textbf{Vercel} for hosting and continuous deployment.
\end{itemize}

\subsection{Expected Benefits}

Relative to the spreadsheet workflow it replaces, GreenStep is expected to
deliver faster action logging, gamified engagement that sustains
participation, live leaderboards that create healthy competition, one-click
administrative CSV export, and very low maintenance overhead thanks to its
fully managed, serverless infrastructure.

% ─────────────────────────────────────────────────────────────────────────────
\section{Current Functionality \& Features}
% ─────────────────────────────────────────────────────────────────────────────

This section describes the features that are implemented and working in the
current build. Each major screen is shown as a captioned screenshot.

\subsection{Authentication and Onboarding}

Authentication is handled by Clerk. New users sign up, sign in, and complete a
first-run onboarding flow that records their department and enrolls them in the
default ``global'' challenge. The landing page (Figure~\ref{fig:landing})
introduces the product and routes signed-in users straight to their dashboard.

\appshot{screenshot-landing.png}{The GreenStep landing page, introducing the
product and providing the sign-in entry point.}{fig:landing}

\subsection{Action Logging}

Authenticated users log sustainable actions across five categories
---Transport, Water, Energy, Recycling, and Food. Each action carries a point
value, and users may optionally attach a note or a photo for a small bonus.
Users browse and search a catalog of pre-defined actions
(Figure~\ref{fig:logaction}) and then confirm a specific action on its own
page.

\appshot{screenshot-log-action.png}{The action-logging screen, where users
browse the five categories and select a specific action to log.}{fig:logaction}

\subsection{Dashboard, Streaks, and Badges}

The dashboard (Figure~\ref{fig:dashboard}) is the home screen. It shows the
user's level and progress, current streak, recent personal activity, a
community activity feed, and a performance chart comparing the user against the
group average (rendered with Recharts). Daily engagement builds a
\textbf{streak}, and users earn \textbf{badges} for milestones such as their
first action, point thresholds, category-specific counts, and streak lengths.

\appshot{screenshot-dashboard.png}{The main dashboard showing level progress,
streak, recent activity, and the performance chart.}{fig:dashboard}

\subsection{Teams and Admin Approval}

Within a challenge, users join existing teams or submit a request to create a
new team. Team-creation requests enter a \texttt{PENDING} state and must be
approved (or rejected) by an administrator, giving organizers governance over
team structure.

\subsection{Leaderboards}

GreenStep provides both \textbf{team} and \textbf{individual} leaderboards
(Figure~\ref{fig:leaderboard}). Individual rankings offer several views
---all-time points, rising performers, longest streaks, and category
diversity---and both leaderboards can be filtered by category and time range.

\appshot{screenshot-leaderboard.png}{The leaderboard screen, showing team and
individual rankings with category and time-range filters.}{fig:leaderboard}

\subsection{Challenges and Join-by-Code}

Every user is auto-enrolled in a permanent global challenge and may join
additional challenges using a unique join code. The active challenge scopes
the dashboard, leaderboards, and activity feeds, and users can switch between
challenges they belong to.

\subsection{Administrator Tools}

Administrators---identified by a Clerk \texttt{publicMetadata} role---have
access to an admin screen (Figure~\ref{fig:admin}) for creating challenges,
reviewing and approving team requests, and exporting participation data to a
CSV file for external reporting.

\appshot{screenshot-admin.png}{The administrator screen for creating
challenges, reviewing team requests, and exporting CSV data.}{fig:admin}

% ─────────────────────────────────────────────────────────────────────────────
\section{To-Do List \& Future Enhancements}
% ─────────────────────────────────────────────────────────────────────────────

The following items were out of scope for this semester or surfaced during
development. They are grouped by aspect.

\subsection{Technical}

\begin{itemize}
    \item \textbf{Automated testing.} The project currently relies on manual
    testing. Adding unit and end-to-end tests (e.g.\ for authentication, action
    logging, and leaderboard aggregation) and wiring them into CI would catch
    regressions automatically.
    \item \textbf{Further dashboard decomposition.} The main
    \texttt{DashboardClient.tsx} component was recently reduced (from roughly
    2,300 to about 2,080 lines) by extracting its types, helpers, and the
    \texttt{useCountUp} hook into separate modules. The next step is to split
    it into per-tab components (Home, Leaderboard, Activity, Teams).
    \item \textbf{Database migrations.} The schema is currently applied with
    \texttt{prisma db push}. Moving to versioned \texttt{prisma migrate}
    migrations would provide a reproducible schema history for production.
\end{itemize}

\subsection{Functional}

\begin{itemize}
    \item \textbf{Undo / edit a logged action.} Users cannot currently correct
    or remove an action after logging it.
    \item \textbf{Richer admin dashboard.} Expand the admin tools with
    participation analytics and per-team summaries.
    \item \textbf{Real challenge scheduling.} Honor the \texttt{startDate} and
    \texttt{endDate} fields to open and close challenges automatically.
    \item \textbf{User-selectable action dates.} Allow back-dating an action;
    actions are currently timestamped automatically.
\end{itemize}

\subsection{Usability}

\begin{itemize}
    \item \textbf{Reminders and notifications.} Nudge users to maintain their
    streaks and re-engage lapsed participants.
    \item \textbf{Accessibility audit.} Conduct a formal review against WCAG
    guidelines (keyboard navigation, color contrast, screen-reader labels).
    \item \textbf{Environmental-impact feedback.} Show simple, category-based
    impact facts when an action is logged.
\end{itemize}

\subsection{Future Expansion}

\begin{itemize}
    \item \textbf{Email digests.} Periodic summaries of team standings and
    personal progress.
    \item \textbf{Deeper analytics.} Trend analysis and reporting for
    organizers across challenges.
    \item \textbf{Multi-organization scaling.} Offer GreenStep as a
    multi-tenant service so other organizations can run their own challenges.
\end{itemize}

% ─────────────────────────────────────────────────────────────────────────────
\begin{thebibliography}{9}
% ─────────────────────────────────────────────────────────────────────────────

\bibitem{nextjs}
Vercel, ``Next.js Documentation,'' 2024. [Online]. Available:
\url{https://nextjs.org/docs}

\bibitem{clerk}
Clerk, ``Clerk Documentation,'' 2024. [Online]. Available:
\url{https://clerk.com/docs}

\bibitem{prisma}
Prisma, ``Prisma ORM Documentation,'' 2024. [Online]. Available:
\url{https://www.prisma.io/docs}

\bibitem{tailwind}
Tailwind Labs, ``Tailwind CSS Documentation,'' 2024. [Online]. Available:
\url{https://tailwindcss.com/docs}

\bibitem{recharts}
Recharts, ``Recharts Documentation,'' 2024. [Online]. Available:
\url{https://recharts.org/}

\bibitem{neon}
Neon, ``Neon Serverless Postgres Documentation,'' 2024. [Online]. Available:
\url{https://neon.tech/docs}

\bibitem{cisc489}
E. Adekoya, K. Alkhaleefa, and J. Dunlap, ``Employee Sustainability Challenge
App: Final Project Documentation,'' CISC489 Topics in Human-Computer
Interaction, University of St.\ Thomas, Dec.\ 2025.

\end{thebibliography}

\end{document}
